\section{Compound Statements}\label{sec:CompoundStatements}
Compound statements\index{compound statement} are statements that contain lists of statements enclosed in curly braces \lstinline|{ <statement, … > }|. A compound statement (also called a \bdq{code block}\index{code block|see {compound statement}} or \bdq{block statement}\index{block statement| see {compound statement}}) typically appears as the body of another statement, such as the \lstinline|if| or \lstinline|for| statements, but it may appear also on it own\footnote{In Java, a block also indicates a scope\index{scope}; a local variable lives only in the scope it was declared in, and all inner scopes contained in it.}.

The block's opening curly brace \bdq{\{} appears alone at a line by itself with the same indentation as the statement it belongs to and the corresponding closing curly brace \bdq{\}} is placed at a new line with the same indentation. All code inside such a block is indented one level more than the curly braces.

Large compound statements (long code blocks) should be marked with a label\index{label} and should get an end comment, like in the examples below: 
\begin{lstlisting}
ForeverLoop: while( true )
{
    … // Lots of lines
}   //  ForeverLoop:

CountrySelector: switch( countryCode )
{
    case DE: …
    case EN: …
    … // All the other countries
    default: …
}   //  CountrySelector:

ProcessLoop: for( final var entry : entries )
{
    … // Lots of lines
}   //  ProcessLoop:
\end{lstlisting}

You may have seen something like this in existing code:
\begin{lstlisting}
// DISCOURAGED!
if( thisConditionIsTrue )
{
    // Lots of lines
    …
}   //	if( thisConditionIsTrue ) 
\end{lstlisting}
As long as the \lstinline|if| clause (or whatever statement starts the block) will not change, this pattern serves the same purpose as my variant with the label. But when you change the \lstinline|if|, you have to modify the end comment accordingly – something that is forgotten much too often!

As there is rarely a reason to change the label (the purpose of the block remains the same, no matter what the condition will be), it will not happen that the end comment will not match the start \bsq{comment} of the block.

See the chapters \tqfullref{sec:LabelsAndBreakStatements}, \tqfullref{sec:TrailingOrEndOfLineComments}, and \tqfullref{sec:CommentsWhen} for some additional details. 

For more examples regarding the formatting of compound statements refer in particular to the chapters \tqfullref{sec:IfStatements}, \tqfullref{sec:SwitchStatements}, \tqfullref{sec:ForStatements}, and \tqfullref{sec:WhileStatements}.

Empty compound statements must be avoided. At least a comment must be provided:
\begin{lstlisting}
{
    /* Nothing to do here */
}
\end{lstlisting}

%------------------------------------------------------------------------------

\subsection{\bsq{if}, \bsq{if-else}, \bsq{if-else~if-else} Statements}\label{sec:IfStatements}

We already discussed the use of blanks, round brackets, curly braces and linefeeds in \tqfullvref{sec:WhiteSpace}.

\keeplisting{11}
The base form for the \lstinline|if| statement is like this:
\begin{lstlisting}
if( <condition>)
{
    // <condition> is true
    <statements>;
}
else
{
    // <condition> is false
    <statements>;
}
\end{lstlisting}

\keeplisting{7}
If there is nothing to do in case the condition results as false, the \lstinline|else| clause must be omitted:
\begin{lstlisting}
if( <condition>)
{
    // <condition> is true
    <statements>;
}
\end{lstlisting}

\keeplisting{8}
Negate the condition if there is nothing to do when the original condition is true; do not let the first statement block empty:
\begin{lstlisting}
// RECOMMENDED!!
if( !<condition>)
{
    // <condition> is false; that means !<condition> is true
    <statements>;
}
\end{lstlisting}

\keeplisting{12}
Instead of:
\begin{lstlisting}
// AVOID!!
if( <condition>)
{
    // <condition> is true
    // Nothing to do here!
}
else
{
    // <condition> is false
    <statements>;
}
\end{lstlisting}

\keeplisting{3}
If there is no \lstinline|else| clause, and only a single statement, you can write the \lstinline|if| statement like this:
\begin{lstlisting}
if( <condition> ) <statements>;
\end{lstlisting}

The curly braces are mandatory if it is not possible to write the whole \lstinline|if| statement within a single line.

\keeplisting{22}
In case there are more conditions to check, but a \lstinline|switch| statement cannot be used due to the data types involved or the logic for the conditions, use this form:
\begin{lstlisting}
if( <condition1>)
{
    <statements>;
}
else if( <condition2>)
{
    <statements>;
}
else if( <condition3>)
{
    <statements>;
}
else if( <condition#>)
{
    <statements>;
}
else
{
    <statements>;
}
\end{lstlisting}

Of course you can repeat the \lstinline|else if| as often as you need, but at some point the readability suffers seriously, and you should think about another logic. Also the final \lstinline|else| clause is not mandatory.

\keeplisting{33}
It is not recommended to write \verb#if-else# like this, although it seems to be more conformant to the basic rule:
\begin{lstlisting}
// NOT RECOMMENDED!!
if( <condition1>)
{
    <statements>;
}
else 
{
    if( <condition2>)
    {
        <statements>;
    }
    else 
    {
        if( <condition3>)
        {
            <statements>;
        }
        else 
        {
            if( <condition#>)
            {
                <statements>;
            }
            else
            {
                <statements>;
            }
        }
    }
}        
\end{lstlisting}

The increasing indendation makes this even worse to read than the recommended variant, particularly when the code lines get longer.

\subsubsection{Principles}
\begin{enumerate}[label=P\arabic*.]
\item{Curly braces are mandatory when the \lstinline|if| statement has more than one line. They are always required for the \lstinline|else| clause.}
\item{No empty \lstinline|if| clause; negate the condition if necessary.}
\item{No blanks between the \lstinline|if|\index{if} keyword and the opening round bracket. 

Refer to \tqfullvref{sec:ParameterParenthesis} on how to place the white space.}
\item{The opening curly brace for the body statement of the \lstinline|if| and the optional \lstinline|else| clause is on the same indentation level as the keyword itself. In case of the \verb#if-else# construct, on the level of the \lstinline|else|.}
\item{The statements itself are indented one level relative to the curly braces.}
\end{enumerate}

Please refer also to chapter \tqvref{sec:TheTernaryOperator} that discusses the ternary \verb#?# operator.

%==============================================================================
% End of file!!
%==============================================================================

\subsection{\bsq{switch} Statements}\label{sec:SwitchStatements}
This chapter describes how to format a \lstinline|switch|\index{switch} statement together with the related \lstinline|case|\index{case} and \lstinline|default|\index{switch!default} statements.

Beginning with the first preview in Java~12, an alternative syntax for \lstinline|switch|\index{switch} was introduced, so that we have now two (in fact, three) different forms of that construct.

Some sources (e.g. \autocite{ORACLE_DOC_SWITCHEXPRESSIONS}) refer to the different syntax variants by the form of the \lstinline|switch|\index{switch} labels as ``\verb#case L:#'' (\bsq{old}) and ``\verb#case L ->#'' (\bsq{new}). 

\begin{tabular}{ll}
\rule[-1ex]{0pt}{2.5ex} ``\verb#case L:#'' or \bsq{old} \lstinline|switch| & ``\verb#case L ->#'' or \bsq{new} \lstinline|switch| \\ 
\rule[-1ex]{0pt}{2.5ex}
\begin{minipage}{0.48\textwidth}
\begin{lstlisting}[numbers=left]
switch( <selector> )
{
    case <label1>:
        <statements1>; 
        break;
    
    case <label2a>:
    case <label2b>: 
        <statements2>; 
        break;
    
    default: <statements>; break;
}
\end{lstlisting}
\end{minipage} & 

\begin{minipage}{0.48\textwidth}
\begin{lstlisting}
switch( <selector> )
{
    case <label1> -> <statements1>;
    
    
        
    case <label2a>, <label2b> 
        -> <statements2>;


    
    default -> <statments>
}
\end{lstlisting}
\end{minipage} \\ 
\end{tabular} 

%==============================================================================
% End of file!!
%==============================================================================


I will discuss in chapter \tqfullvref{sec:TheSwitchStatement} how and when to use which form.

%------------------------------------------------------------------------------

\subsubsection{The traditional Form of \bsq{switch}}
Originally, a \lstinline|switch|\index{switch}\index{switch!statement} statement in \Java looked like this (the ``\verb#case L:#'' variant):
\begin{lstlisting}[numbers=left]
switch( <selector> )
{
    case <switchlabel1>:
        <statements1>;
        // Falls through!

    case <switchlabel2>:
        <statements2>;
        break;

    case <switchlabel3>: <singleStatement>; break;

    case <switchlabel4>:
    {
        <statements>;
        break;
    }

    case <switchlabel5>: // Also a fall-through
    case <switchlabel6>:
    {
        <statements>;
        break;
    }

    default:
       <statements>;
       break;
}
\end{lstlisting}

First, these coding conventions make the blank line above each \lstinline|case|\index{case} (except the very first one) and above \lstinline|default| mandatory!
  
If a \lstinline|case| falls through, like in lines~3 to 9, a comment \bdq{\lstinline|// Falls through!|} (or something similar\footnote{Linting tools as well as most IDEs issue a warning in case of a fall-through. Usually, these could be switched off by adding a comment or an annotation to the location of the fall-through. This would be a sufficient replacement for the comment.}) as in line~5 of the sample is required! But basically, such constructs should be avoided when possible, because it also forces that the sequence of the \lstinline|case| clauses cannot be changed, and that fact requires an additional comment to the \lstinline|switch| itself. In fact here you should consider to ignore the \bdq{DRY~Principle}\index{DRY Principle} – we will see this later – and repeat the code for the second \lstinline|case| on the first one:
\begin{lstlisting}
// BETTER
switch( <selector> )
{
    case <switchlabel1>:
        <statements1>;
        <statements2>;
        break;

    case <switchlabel2>:
        <statements2>;
        break;

    default: …
}
\end{lstlisting}
Or you put the code that is common for both cases into a \lstinline|private| method and call that from both branches. 

Different to the case shown in lines~3 to 9, no comment is required in the case shown in lines~19 and 20: here we have two \lstinline|case| clauses that triggers \textit{exactly the same} action, and we can even omit the blank line between them (and also the comment from line~19).

For longer \lstinline|switch| statements it is highly recommended to use a label with the \lstinline|break| statement\footnote{In fact, I recommend to always use a label with \lstinline|break|; see chapter \tqfullref{sec:LabelsAndBreakStatements} on this topic.}; this can look like this:
\begin{lstlisting}
final Color color;
ColorSelector: switch( colorIndex )
{
    case 1: color = RED; 
        break ColorSelector;
        
    case 2: color = BLUE; 
        break ColorSelector;
        
    case 3: color = YELLOW; 
        break ColorSelector;
        
    case 4: color = GREEN; 
        break ColorSelector;
        
    default: color = NONE; 
        break ColorSelector;
}   //  ColorSelector:
\end{lstlisting}

The \lstinline|break| in the \lstinline|default| branch is redundant, but it prevents a fall-through error if later another \lstinline|case| branch is added – although it should be assured that the \lstinline|default| branch is always the last branch in the \lstinline|switch| statement. Of course no \lstinline|break| is required when the branch is left by throwing an exception.\footnote{A \lstinline|break| would be also obsolete when the branch is left through a \lstinline|return| statement, but in that case, the \lstinline|return| would not be the last statement in the method, and this is another requirement from this coding conventions!}

%==============================================================================
% End of file!!
%==============================================================================

\subsubsection{The new Form of \bsq{switch}}\label{sec:NewSwitch}
The alternative syntax for \lstinline|switch|\index{switch} comes in two flavours (as \textit{statement}\index{switch!statement} and \textit{expression}\index{switch!expression}).

The new \lstinline|switch|\index{switch!statement} statement looks like this:

\begin{lstlisting}[numbers=left]
// switch statement
switch( <selector> )
{
    case <switchlabel1> -> <singleStatement>;

    case <switchlabel2>, <switchlabel3> -> <singleStatement>;

    case <switchlabel4> ->
    {
        <statements>;
    }

    case <switchlabel4>, <switchlabel5> ->
    {
        <statements>;
    }

    default -> <singleStatement>;
}
\end{lstlisting}

And this is the new \lstinline|switch|\index{switch!expression} expression:

\begin{lstlisting}[numbers=left]
// switch expression
var result = switch( <selector> )
{
    case <switchlabel1> -> <expression>;

    case <switchlabel2>, <switchlabel3> -> <expression>;

    case <switchlabel4> ->
    {
        <statements>;
        yield <value>;
    }

    case <switchlabel4>, <switchlabel5> ->
    {
        <statements>;
        yield <value>;
    }

    default -> <expression>;
}
\end{lstlisting}

Of course the \lstinline|default| in line~18 can be followed by a statement block as the \lstinline|case| in line~8, and the \lstinline|default| in line~40 can have a complex expression like the \lstinline|case| in line~28.

As this format does not allow a fall-through in cases where the branches for two or more labels are doing the same stuff, it is possible here to have more than one switch label per \lstinline|case|, as shown in lines~6, 13, 26 and 34.

If used with pattern matching (refer to \autocite{ORACLE_DOC_PATTERNMATCHING}), a \lstinline|switch| looks like this:
\begin{lstlisting}[numbers=left]
var selector = <anObject>
// switch statement
switch( selector )
{
    case null -> <singleStatement>;
    
    case <class1> <name1> -> <singleStatement>;

    case <class2> <name2> ->
    {
        <statements>;
    }

    default -> <singleStatement>;
}

// switch statement
var result = switch( selector )
{
    case null -> <expression>;

    case <class1> <name1> -> <expression>;

    case <class2> <name2> ->
    {
        <expressions>;
        yield <value>;
    }

    default -> <expression>;
}
\end{lstlisting}

Same as for the \lstinline|default| branch, the \lstinline|case null| branch can be a statement block or a complex expression.

And an expression can also be always a \lstinline|throw|.


%==============================================================================
% End of file!!
%==============================================================================

\subsubsection{\bsq{switch} and Pattern Matching}\label{sec:PatternMatchingSwitch}
Pattern matching\index{pattern matching} (refer to \autocite{ORACLE_DOC_PATTERNMATCHING}\index{pattern matching}) is covered in more detail in chapter \tqfullvref{sec:PatternMatching}. Here I want to discuss only the basic formatting of the \lstinline|switch|\index{switch} when used with pattern matching.

\keeplisting{14}
A \lstinline|switch| statement\index{switch!statement} for the ``\verb#case L:#'' variant would look like this:
\index{switch!case}\index{switch!default}\index{switch!break}
\begin{lstlisting}
var selector = <anObject>
switch( selector )
{
    case <class1> <name1>: <singleStatement>;

    case <class2> <name2>:
    {
        <statements>;
    }

    default: <singleStatement>;
}
\end{lstlisting}

\keeplisting{15}
For the ``\verb#case L:#'' variant, a \lstinline|switch| statement\index{switch!statement} is written as below:
\index{switch!case}\index{switch!default}
\begin{lstlisting}
var selector = <anObject>
switch( selector )
{
    case null -> <singleStatement>;
    
    case <class1> <name1> -> <singleStatement>;

    case <class2> <name2> ->
    {
        <statements>;
    }

    default -> <singleStatement>;
}
\end{lstlisting}

\keeplisting{16}
And an ``\verb#case L->#'' variant \lstinline|switch| expression\index{switch!expression} would have this form:
\index{switch!case}\index{switch!default}\index{switch!yield}
\begin{lstlisting}
var selector = <anObject>
var result = switch( selector )
{
    case null -> <expression>;

    case <class1> <name1> -> <expression>;

    case <class2> <name2> ->
    {
        <expressions>;
        yield <value>;
    }

    default -> <expression>;
}
\end{lstlisting}

Same as for the \lstinline|default|\index{switch!default} branch, the \lstinline|case null| \index{switch!case null} branch can be a statement block or a complex expression.

And an expression can also be always a \lstinline|throw|\index{throw} statement.



%==============================================================================
% End of file!!
%==============================================================================

\subsubsection{General formatting Rules for \bsq{switch}}

The general use of blanks, round brackets, curly braces and linefeeds was already discussed in \tqfullvref{sec:WhiteSpace}. The curly braces after the initial \lstinline|switch|\index{switch} are required by the Java syntax and not optional. In case of a \lstinline|switch| statement\index{switch!statement} the opening and closing curly braces are on the same indentation level than the \lstinline|switch| keyword, while for an expression\index{switch!expression} they are aligned with the begin of the respective statement. In both cases, the curly braces are on lines of their own.

The branches (\lstinline|case|\index{switch!case} or \lstinline|default|\index{switch!default}) are indented one level and will be separated from each other by a blank line, to improve the readability.

Although not required (at least not for the ``\verb#case L:#'' variant), curly braces should be used for complex branch statements/expressions, also for readbility. The curly braces for the branches are indented to the same level to the same level as the \lstinline|case|/\lstinline|default|\index{switch!case}\index{switch!default} keywords, while the branch code itself is indented one additional level.

The code for a branch should be short; I suggest that a branch should be limited to less than ten lines. If it is just onle line and it is short enough, the whole branch should be written as a single line, in particular for the ``\verb#case L ->#'' variant:

\index{switch!case}\index{switch!break}
\begin{lstlisting}
case <switchlabel>: <singleStatement>; break;
\end{lstlisting}

\index{switch!case}
\begin{lstlisting}
case <switchlabe> -> <singleStatement>;
\end{lstlisting}

\index{switch!case}
\begin{lstlisting}
case <switchLabel> -> <expression>;
\end{lstlisting}

If necessary, create a private method with the branch code and call that from the branch.

The branches should be ordered somehow logically (alphabetically, by ordinal number,~…), and an existing \lstinline|case null|\index{switch!case null} branch should be the first, while the \lstinline|default|\index{switch!default} branch should be the last one. This makes it easier to apply changes to the \lstinline|switch|\index{switch} later.

\subsubsection{Principles}
\begin{enumerate}[label=P\arabic*.]
\item{No blanks between the \lstinline|switch|\index{switch} and the opening round bracket. Refer to \tqfullvref{sec:ParameterParenthesis} on how to place the white space.}
\item{Curly braces are recommended when the branch statement has more than one line.}
\item{The curly braces for a branch statement are on the same indentation level as the keyword itself.

The branch statements itself are indented one level relative to the curly braces.}
\item{The code for a branch should be short.}
\item{Apply a logical order to the branches.

Place \lstinline|case null|\index{switch!case null} as first, and \lstinline|default|\index{switch!default} as last branches.}
\item{The blank line separating the branches is mandatory.}
\item{No empty branch! A branch must contain at least a comment.

An exception for the ``\verb#case L:#'' variant are branches with the same code.}
\item{Avoid ‘fall-through’!}
\item{Although it looks like one, a \lstinline|switch| expression\index{switch!expression} is not a lambda\index{lambda}.

That means that local variables that are referenced in a \lstinline|switch| expression do not have to be effectively final.}
\end{enumerate}

%==============================================================================
% End of file!!
%==============================================================================


%==============================================================================
% End of file!!
%==============================================================================

\subsection{‘for’ Statements}\label{sec:ForStatements}
Since Java~5 the language knows two distinct types of \lstinline|for|\index{for} loops, both described in \autocite{ORACLE_DOC_LANGUAGE_SPECIFICATION:For}.

The basic \lstinline|for| statement\index{for!basic} looks like this:
\begin{lstlisting}[numbers=left]
for( <initialization>; <condition>; <update> )
{
    <statements>;
}

for( <initialization>; <condition>; <update> ) <singleStatement>;
\end{lstlisting}
For the details refer to \autocite{ORACLE_DOC_LANGUAGE_SPECIFICATION:ForBasic}.

An enhanced \lstinline|for| statement\index{for!enhanced} (defined in \autocite{ORACLE_DOC_LANGUAGE_SPECIFICATION:ForEnhanced}) would be written like below:

\begin{lstlisting}[numbers=left]
for( <declaration> : <iterable> )
{
    <statements>;
}

for( <declaration> : <iterable> ) <singleStatement>;
\end{lstlisting}

The use of blanks, round brackets, curly braces and linefeeds was already discussed in \tqfullvref{sec:WhiteSpace}.

Curly braces can be omitted only when there is just a single simple statement to be executed in the loop. In addition, this single statement has to be written completely into the same line as the \lstinline|for| itself. In any other case, the curly braces are mandatory.

An empty loop body can be written like this:
\begin{lstlisting}[numbers=left]
for( <initialization>; <condition>; <update> )
{
    // Deliberately empty
}

for( <initialization>; <condition>; <update> );

for( <declaration> : <iterable> )
{
    // Deliberately empty
}

for( <declaration> : <iterable> );
\end{lstlisting}

In case the longer forms are used, the comment is mandatory, because empty code blocks are forbidden. I recommend the shorter forms (lines~6 and 13).

When explicitly using an iterator\index{iterator}\index{java.util!Iterator}\autocite{ORACLE_DOC_ITERATOR_INTERFACE} with the (basic) \lstinline|for|\index{for!basic} loop, this should look like this:
\begin{lstlisting}
for( final var i = source.iterator(); i.hasNext(); )
{
    …
}
\end{lstlisting}

When using the comma operator in the initialization or update clause of a \lstinline|for|\index{for} statement, avoid the complexity of using more than three variables. If needed, use separate statements before the \lstinline|for|-loop (for the initialization clause) or at the end of the loop (for the update clause).

Examples:
\begin{lstlisting}
for( var i = 0, j = target.length; i < source.length && j >= 0; ++i, --j )
{
    target [j] = source [i];
}

var proceed = true;
var i = source.iterator();
for( var j = 0; proceed; )
{
    target [j++] = i.next();
    proceed = i.hasNext() && j < target.length;
}
\end{lstlisting}

\subsubsection{Principles}
\begin{enumerate}[label=P\arabic*.]
\item{one}
\item{two}
\item{three}
\end{enumerate}

\lstinline|for| loops are discussed in more detail in \tqfullvref{sec:TheForStatement}.

Examples:
\begin{lstlisting}[numbers=left]
// Ok
for( var i = 0; i < max; ++i ) sum += i;
for( final var s : texts ) System.out.println( s );

// DISCOURAGED
for( var i = 0; i < max; ++i ) if( v [i] > 0 ) sum += v [i];
for( final var s : texts ) if( !s.empty() ) out.println( s );

// AVOID!!!
for( var i = 0; i < max; ++i ) if( v [i] > 0 )
{
    sum += v [i];
}
else
{
    sum -= v [i];
}

for( final var s : texts ) if( !s.empty() )
{
    System.out.print( "Value: " );
    System.out.println( s );
}
\end{lstlisting}
The samples in lines~6 and 7 are not correct according to this rule, because the \lstinline|if| statement is not a simple statement..

An empty \lstinline|for|-loop (one in which all the work is done in the initialization, condition, and update clauses) should have the following form:\footnote{An \textit{enhanced} \lstinline|for|-loop without body does not make that much sense (that mentioned above is a classical \lstinline|for|-loop). It may be possible to write an implementation of \lstinline|java.lang.Iterable| with an iterator that does something as a side effect of \lstinline|hasNext()| or \lstinline|next()|, but this would break the contract of the interface \lstinline|java.lang.Iterator|.}
\begin{lstlisting}
for( <initialization>; <condition>; <update> );
\end{lstlisting}


%==============================================================================
% End of file!!
%==============================================================================

\subsection{\bsq{while} and \bsq{do} Statements}\label{sec:WhileStatements}
The \lstinline|do|\index{while!do} statement does never come without a trailing \lstinline|while|\index{while} statement; this is defined in \autocite{ORACLE_DOC_LANGUAGE_SPECIFICATION:WhileDo}. The statement \lstinline|while|\index{while} itself is defined in \autocite{ORACLE_DOC_LANGUAGE_SPECIFICATION:While}. Both are loop statements, the difference is only where the termination condition is checked.

With \lstinline|do|\index{while!do} you define an exit-controlled loop\index{exit-controlled loop}\footnote{In German: \bdq{nicht-abweisende Schleife}.}; this means that the loop body is executed at least once:
\begin{lstlisting}
do
{
    <statements>;
}
while( <condition> );
\end{lstlisting}
Although writing
\begin{lstlisting}
do <statement>; while( <condition> );
\end{lstlisting}
is syntactically correct (note the semicolon after \verb#<statement># and before \lstinline|while|) you should always surround the loop body with curly braces when using \lstinline|do|\index{while!do}.

The entry-controlled \lstinline|while| loop\index{entry-controlled loop}\index{while}\footnote{In German: \bdq{abweisende Schleife}.} is used more often than the exit-controlled form; it looks like this:
\begin{lstlisting}
while( <condition> )
{
    <statements>;
}

while( <condition> ) <singleStatement>;
\end{lstlisting}

Same as for the \lstinline|for| statement (see \hyperref[sec:ForStatements]{there}), \bdq{single statement} also means \bdq{simple statement}. If the loop body is more complex, curly braces are mandatory; they are placed to lines of their own, indented to same level as the keyword, while the body itself has to be indented one additional level. The curly braces can be omitted only for a body consisting of a single statement that is written to the same line as the \lstinline|while|\index{while} statement.

The form
\begin{lstlisting}
// DISCOURAGED
while( <condition> )
    <statements>;
\end{lstlisting}
is syntactically correct, but have to be avoided. It is error prone in case the code will be adjusted at a later time.

The general use of blanks, round brackets, curly braces and linefeeds was already discussed in \tqfullvref{sec:WhiteSpace}.

When the loop body is larger, you should consider to use a label\index{label}, also when you use \lstinline|break|\index{break} or \lstinline|continue|\index{continue} inside the loop body:
\begin{lstlisting}[numbers=left]
<Label>: while( <condition> )
{
    if( <condition1> ) continue <Label>;
    if( <condition2> ) break <Label>;
    <statements>
}   //  <Label>:

<Label>: do
{
    if( <condition1> ) continue <Label>;
    if( <condition2> ) break <Label>;
    <statements>
} 
while( <condition> );
//  <label>
\end{lstlisting}

The body of a \lstinline|while| loop\index{while}\index{while!do} must not be empty! Or, more precisely, it must contain executable code! Usually a comment is sufficient to document that nothing happens inside that code block, but for a \lstinline|while|\index{while!empty loop} loop, this is not sufficient. The details are explained in the chapter \tqfullvref{sec:EmptyWhile}.

\subsubsection{Principles}
\begin{enumerate}[label=P\arabic*.]
\item{No blanks between the \lstinline|for|\index{for} and the opening round bracket. The blanks after the semi-colons are mandatory.

Refer to \tqfullvref{sec:ParameterParenthesis} on how to place the white space.}
\item{If the loop body is a simple single statement, it can be written without curly braces to the same line as the \lstinline|for|\index{for} itself.

Otherwise curly braces are mandatory.}
\item{The curly braces are on lines of their own, with the same indentation level as the \lstinline|for|\index{for} keyword. The code of the body will be indented one level.}
\item{Do not use a \lstinline|for|\index{for} loop to consume CPU cycles only!}
\item{Empty \lstinline|for|\index{for!empty} loops do all the work in the initialisation, condition and update clauses.

Prefer a \lstinline|while|\index{while} loop over an empty \lstinline|for|\index{for!empty} loop.}
\item{Keep the \lstinline|for|\index{for} statement simple.

Avoid using the comma operator in the initialisation and update clause.}
\item{Keep the loop body short.

Use a label\index{label} for longer loops.}
\item{Use the \lstinline|for|\index{for} loop variant that fits best to your current problem.

The enhanced \lstinline|for|\index{for!enhanced} is not generally better than the basic variant, it just matches better with some iteration tasks.}
\end{enumerate}

\lstinline|for| loops are discussed in more detail in chapter \tqvref{sec:TheForStatement}.










\subsection{‘do-while’ Statements}
A \lstinline|do-while| statement should look like below:







A \lstinline|while| statement should have one of the following forms:
\begin{lstlisting}[numbers=left]
while( <condition> )
{
    <statements>;
}
while( <condition> ) <singleStatement>
\end{lstlisting}

Again, the curly braces can only be omitted in case there is only a single statement, written on the same line than the \lstinline|while| statement that has to be executed in the loop.

An empty \lstinline|while|-loop\footnote{An empty \lstinline|while|-loop usually makes no sense as the compiler optimises it away. Only when the condition causes side effects, it will be executed. But such an implementation is difficult to understand and should be avoided.} should have the following form: 
\begin{lstlisting}
while( <condition> );
\end{lstlisting}

Longer \lstinline|while|-loops should use a label:
\begin{lstlisting}
LoopLabel: while( proceed )
{
    // Lots of code lines here …
}   //  LoopLabel:
\end{lstlisting}

Chapter \tqvref{sec:TheWhileStatement} has some more about \lstinline|while| loops.

%==============================================================================
% End of file!!
%==============================================================================

\subsection{\bsq{throw}, \bsq{try-catch} and \bsq{try-catch-finally} Statements}\label{sec:TheTryCatchStatements}
Throwing an exception and catching it somewhere else in the code is the foundation of error handling in the \Java programming language. The \verb#try-catch-finally#\index{try}\index{try!catch}\index{try!finally} statement is the core element here, together with the \lstinline|throw|\index{throw} statement and the exceptions itself. Error handling is covered in detail in chapter \tqfullvref{sec:ErrorHandling}.

The \lstinline|try|\index{try} statement is defined in \autocite{ORACLE_DOC_LANGUAGE_SPECIFICATION:Try}, and the definitions for the \lstinline|throw|\index{throw} statement can be found in \autocite{ORACLE_DOC_LANGUAGE_SPECIFICATION:Throw}.

\subsubsection{Formatting \bsq{throw}}
There is not much to say about the \lstinline|throw|\index{throw} statement in regards of formatting:
\begin{lstlisting}
throw <exception>;
\end{lstlisting}

That's it!

\subsubsection{Formatting \bsq{try-catch-finally}}
A \verb#try-catch-finally#\index{try}\index{try!catch}\index{try!finally} statement has the following format:
\begin{lstlisting}
try
{
    <statements>;
}
catch( final <ExceptionClass> e )
{
    <statements>;
}
finally
{
    <statements>
}
\end{lstlisting}
There can be more than one \lstinline|catch|\index{try!catch} clauses, but just one \lstinline|finally|\index{try!finally} clause. If there is at least one \lstinline|catch|\index{try!catch} clause, the \lstinline|finally|\index{try!finally} clause can be omitted. And \lstinline|catch|\index{try!catch} clauses are optional if there is a \lstinline|finally|\index{try!finally} clause.

A \verb#try-with-resources#\index{try-with-resources}\index{try} looks basically like this:
\begin{lstlisting}
try( <resource reference> )
{
    <statements>;
}
catch( final <ExceptionClass> e )
{
    <statements>;
}
finally
{
    <statements>
}
\end{lstlisting}
Again, there can be more than one \lstinline|catch|\index{try!catch} clauses, but just one \lstinline|finally|\index{try!finally} clause. But for a \verb#try-with-resources#\index{try-with-resources}\index{try}, both \lstinline|catch|\index{try!catch} clause and \lstinline|finally|\index{try!finally} clause can be omitted.

For a detailed description of \verb#try-with-resources#\index{try-with-resources}\index{try} refer to chapter \tqfullvref{sec:TryWithResources}.

The use of blanks, round brackets, curly braces and linefeeds was already discussed in \tqfullvref{sec:WhiteSpace}.

The curly braces around the code blocks for the \lstinline|try|\index{try}, the \lstinline|catch|\index{try!catch}, and the \lstinline|finally|\index{try!finally} are syntactically required. The opening and the closing braces will be placed to lines of their own, indented to the level of the respective keyword. The code inside will be indented one additional level, as for any other code block.

In case more than one exception type should be handled the same way, these can be combined into one \lstinline|catch|\index{try!catch} block, like this:
\begin{lstlisting}
try
{
    …
}
catch( final <ExceptionClass1> | <ExceptionClass2> e )
{
    <statements>;
}
\end{lstlisting}
The vertical bar is surrounded by blanks.

An \textit{empty} \lstinline|catch|\index{try!catch} block is unacceptable under all circumstances! Refer to chapter \tqfullvref{sec:GeneralExceptionHandling} how to handle exceptions. See also chapter \tqvref{sec:SingleLineComments} about empty blocks and comments.

Usually, the caught exception is referenced in the \lstinline|catch|\index{try!catch} block. If it is not referenced, you should use the underscore (\bdq{\lstinline|_|}) as the \textit{unnamed exception parameter}\autocite{ORACLE_DOC_LANGUAGE_SPECIFICATION:Declarations}\index{unnamed exception parameter}\footnote{From the \jls: \bdq{If a declaration does not include an identifier, but instead includes the keyword \bdq{\lstinline|_|} (underscore), then the entity cannot be referred to by name. […] A local variable, exception parameter, or lambda parameter that is declared using an underscore is called an unnamed local variable, unnamed exception parameter, or unnamed lambda parameter, respectively.}}. This looks like this:
\begin{lstlisting}
try
{
    …
}
catch( final <ExceptionClass> _ )
{
    <statements>;
}
\end{lstlisting}
Of course this is mandatory if the \lstinline|catch|\index{try!catch} block does not contain any executable code.

A \lstinline|finally| block without any executable code is completely obsolete and must be removed.

\subsubsection{Principles}
\begin{enumerate}[label=P\arabic*.]
\item{No blanks between the \lstinline|try|\index{try-with-resources} and the \lstinline|catch|\index{try!catch} keywords and the opening round brackets.

Refer to \tqfullvref{sec:ParameterParenthesis} on how to place the white space.}
\item{The keywords \lstinline|try|\index{try}, \lstinline|catch|\index{try!catch} and \lstinline|finally|\index{try!finally} are all indented at the same level.

The curly braces  for the respective code blocks are aligned with the keywords, the code is indented one level relative to the curly braces.}
\item{An empty \lstinline|catch|\index{try!catch} block is not allowed!

It must contain at least a comment that explains why the caught exception is not handled.}
\item{A \lstinline|finally|\index{try!finally} block without executable code must be removed.}
\end{enumerate}

%==============================================================================
% End of file!!
%==============================================================================

\subsection{‘synchronized’ Statements}\label{sec:SynchronizedStatements}
The \lstinline|synchronized|\index{synchronized} statement is defined in \autocite{ORACLE_DOC_LANGUAGE_SPECIFICATION:Synchronized}\footnote{The keyword \lstinline|synchronized|\index{synchronized} can be used also as a method modifier; this is explained in \autocite{ORACLE_DOC_LANGUAGE_SPECIFICATION:SynchronizedMethods}.}. A \lstinline|synchronized|\index{synchronized} statement acquires a mutual-exclusion lock on behalf of the executing thread, executes the related code block, then releases the lock. While the executing thread owns the lock, no other thread may acquire the lock – meaning no other thread can enter the block and execute the code in there.

The chapter \tqfullvref{sec:MultiThreading} discusses this topic further.

\keeplisting{5}
Such a \lstinline|synchronized|\index{synchronized} statement may look like this:
\begin{lstlisting}
synchronized( <expression> )
{
    <statements>
}
\end{lstlisting}
The curly braces are mandatory according the Java syntax; their placement follows the rules for any other code block. A single-line \lstinline|synchronized|\index{synchronized} statement is therefore not possible.

\keeplisting{8}
\jls \autocite{ORACLE_DOC_LANGUAGE_SPECIFICATION:Synchronized} says, that the type of \texttt{<expression>} has to be a reference type. That means that
\begin{lstlisting}
// DOES NOT WORK!!
final int i = 5;
synchronized( i )
{
    …
}
\end{lstlisting}
will not work\footnote{In fact, it will cause an error at compile time.}.

\keeplisting{6}
Unfortunately, the compiler will not complain about something like this:
\begin{lstlisting}
// DO NOT DO THIS!!
final StringBuilder sb = new StringBuilder();
synchronized( sb )
{
    …
}
\end{lstlisting}
\keeplisting{7}
And also not about this:
\begin{lstlisting}
// DO NOT DO THIS NEITHER!!
final String s = null;
synchronized( s )
{
    …
}
\end{lstlisting}
Ok, for the second one, you will get a \lstinline|NullPointerException|\index{java.lang!NullPointerException} at runtime, but it may take a lot of time before you spot the error in the first one. But at least most IDEs will issue respective warnings regarding both pattern (if not switched off~…). 

\keeplisting{7}
Even worse to debug would be this pattern, although it could be handy under some circumstances:
\begin{lstlisting}
// TRY TO AVOID!!
synchronized( provideLockObject() )
{
    …
}
\end{lstlisting}
In this case you have to rely on the method \lstinline|provideLockObject()| to return the proper guard object\index{guard object}.

So do not use local or temporary objects with the expression for \lstinline|synchronized|\index{synchronized}.

\keeplisting{16}
In case you declare a field only to provide the monitor for synchronisation, you should name it as \bsq{\texttt{Guard}}\index{guard object}.
\begin{lstlisting}
public final MyClass
{
    private final List<String> m_StringList = new ArrayList<>();
    private static final Object m_StringListGuard = new Object();
    
    public final void myMethod()
    {
        synchronized( m_StringListGuard )
        {
            …
        }
    }   //  myMethod()
}
//  class MyClass
\end{lstlisting}

\subsubsection{Principles}
\begin{enumerate}[label=P\arabic*.]
\item{No blanks between the \lstinline|synchronized|\index{synchronized} keyword and the opening round bracket.

Refer to \tqfullvref{sec:ParameterParenthesis} on how to place the white space.}
\item{The curly braces  for the code block are aligned with the keyword, the code is indented one level relative to the curly braces.}
\item{The \lstinline|synchronised|\index{synchronized} expression should be a final reference field, at best.

If that field exists only for this purpose, it should be named accordingly, as \bsq{\texttt{Guard}}\index{guard object}.}
\item{An empty \lstinline|synchronized|\index{synchronized} block does not make any sense and should be removed.}
\end{enumerate}

%==============================================================================
% End of file!!
%==============================================================================


%==============================================================================
% End of file!!
%==============================================================================
